\chapter{Soft-Margin Support Vector Machine}
\lecture{20}{11 Nov.}{}

\begin{prev}
    Kernel as a shortcut to (transform \(+\) inner product) removes the
    dependence on \(\tilde{d}\) (\autoref{sec:kernel-trick}): it allows a
    spectrum of models, from simple (linear) ones to infinite dimensional
    (Gaussian) ones, all with margin control.
\end{prev}

The kernel trick let us use very powerful transforms, all the way up to the
infinite-dimensional Gaussian one. At the end of the last chapter we saw the
price: with \(\gamma = 100\) the Gaussian SVM classified every training
example correctly, and the boundary clearly overfitted
(\autoref{sec:gaussian-kernel}).

This chapter looks at one cause of that overfitting and removes it. The
hard-margin SVM insists that \emph{every} example lie on its correct side of
the fat boundary. When some labels are noise, the only way to meet this
demand is to bend the boundary around the noisy examples. The
\emph{soft-margin} SVM drops the demand. It allows an example to violate the
margin, and charges a penalty for each violation. A parameter \(C\) sets how
large that penalty is.

The chapter follows the four parts of the deck:
\begin{enumerate}
    \item write down the soft-margin problem (the primal);
    \item derive its dual, which again needs only kernels;
    \item read off what the dual solution tells us: how to compute \(b\),
        and what each \(\alpha_n\) says about its example;
    \item choose the two parameters \((C, \gamma)\) of the Gaussian
        soft-margin SVM.
\end{enumerate}

\paragraph{Roadmap.}
\begin{enumerate}
    \item \textbf{Embedding Numerous Features: Kernel Models}
        \begin{itemize}
            \item[] Lecture 3: Kernel Support Vector Machine ---
                \hlo{\textbf{kernel}} as a shortcut to (transform \(+\) inner
                product) to \hlo{\textbf{remove dependence on
                \(\tilde{d}\)}}: allowing a spectrum of simple
                (\hlo{\textbf{linear}}) models to infinite dimensional
                (\hlo{\textbf{Gaussian}}) ones with margin control
            \item[] \textbf{Lecture 4: Soft-Margin Support Vector Machine}
                --- Motivation and Primal Problem; Dual Problem; Messages
                behind Soft-Margin SVM; Model Selection
        \end{itemize}
    \item Combining Predictive Features: Aggregation Models
    \item Distilling Implicit Features: Extraction Models
\end{enumerate}

\section{Motivation and Primal Problem}\label{sec:soft-primal}

\source{slides p.2--6}

\subsection{Cons of Hard-Margin SVM}

\begin{recall}
    \hlo{\textbf{SVM can still overfit :-(}}
\end{recall}

We first met these two pictures in \autoref{sec:transform-price}. Both
classify the same twenty examples. On the left the boundary is a line
(\(\bm{\Phi}_1\)); on the right it is a fourth-order curve
(\(\bm{\Phi}_4\)).

\begin{center}
    \begin{tabular}{@{}c@{\hspace{1.2em}}c@{\hspace{1.2em}}c@{}}
        \ksvmpanel[3.0]{\softPhiOne}{\softPhiData}{} &
        \begin{minipage}[c]{0.3\textwidth}
            \begin{itemize}
                \item part of reasons: \(\bm{\Phi}\)
                \item other part: \hlp{\textbf{separable}}
            \end{itemize}
        \end{minipage} &
        \ksvmpanel[3.0]{\softPhiFour}{\softPhiData}{}\\[4pt]
        \(\bm{\Phi}_1\) & & \(\bm{\Phi}_4\)
    \end{tabular}
\end{center}

The line gets two examples wrong: the circle in the bottom-left corner and
the cross near the top, just left of the centre. Both look like noise. The
curve gets all twenty right, but only by drawing three separate blue regions
to catch every circle. Most of us would trust the line more. The curve has
overfitted.

The slide names two causes.
\begin{itemize}
    \item \textbf{The transform \(\bm{\Phi}\).} A powerful transform can draw
        very complicated boundaries. The Gaussian kernel of the last chapter
        is even more powerful than \(\bm{\Phi}_4\).
    \item \textbf{Insisting on separable data.} The hard-margin SVM only
        accepts a boundary that puts every example on its correct side. If
        some labels are noise, it has no choice but to bend around those
        examples as well.
\end{itemize}
This chapter deals with the second cause.

\begin{moral}
    If always insisting on \hlp{\textbf{separable}} (\(\Longrightarrow\)
    \hlp{\textbf{shatter}}), have power to \hlo{\textbf{overfit to noise}}.
\end{moral}

Why does `separable' lead to `shatter'? Noise can give the examples any
labels at all. A model that must separate the data whatever the labels are
must be able to produce every possible labelling, and that is exactly what
shattering means. A model with that much power fits noise as easily as it
fits the real pattern, which is overfitting (\autoref{sec:what-overfit}).

\subsection{Give Up on Some Examples}

So we want the SVM to \hlp{\textbf{give up}} on some noisy examples: it
should be allowed to get a few of them wrong when fitting them would cost
too much. We have seen this idea before, in the pocket algorithm
(\autoref{alg:pocket}). Compare the two problems:

\begin{center}
    \begin{tabular}{@{}c@{\qquad}c@{}}
        \textbf{pocket} & \textbf{hard-margin SVM}\\[4pt]
        \(\displaystyle \min_{b, \bm{w}} \sum_{n=1}^N \hlb{\llbracket y_n
            \neq \sign(\bm{w}^\top \bm{z}_n + b) \rrbracket}\)
        &
        \(\begin{aligned}
            \min_{b, \bm{w}} \quad & \green{\tfrac{1}{2} \bm{w}^\top
                \bm{w}}\\
            \text{s.t.} \quad & y_n (\bm{w}^\top \bm{z}_n + b) \ge 1
                \text{ for all } n
        \end{aligned}\)
    \end{tabular}
\end{center}

The pocket algorithm minimizes the number of mistakes. It accepts some
errors, but it does not care about the margin. The hard-margin SVM maximizes
the margin, but it accepts no errors at all. We want something in between,
so we combine the two:
\[
    \begin{aligned}
        \text{combination:} \quad \min_{b, \bm{w}} \quad & \green{\frac{1}{2}
            \bm{w}^\top \bm{w}} + \hlr{C} \cdot \sum_{n=1}^N
            \hlb{\llbracket y_n \neq \sign(\bm{w}^\top \bm{z}_n + b)
            \rrbracket}\\
        \text{s.t.} \quad & y_n (\bm{w}^\top \bm{z}_n + b) \ge 1 \text{ for
            \hlb{\textbf{correct}} } n\\
        & y_n (\bm{w}^\top \bm{z}_n + b) \ge -\infty \text{ for
            \hlb{\textbf{incorrect}} } n
    \end{aligned}
\]
Read it term by term.
\begin{itemize}
    \item The objective has two parts. The green part comes from the
        hard-margin SVM: a smaller \(\bm{w}^\top \bm{w}\) means a larger
        margin. The blue part comes from the pocket: it counts the mistakes,
        and each mistake costs \(C\).
    \item For an example that the hyperplane classifies correctly, we keep
        the usual margin constraint.
    \item For an example that it gets wrong, the constraint becomes \(\ge
        -\infty\), which is always true. In effect that example has no
        constraint any more. This is what `giving up' on it means. We do pay
        for it, though: it adds \(C\) to the objective.
\end{itemize}
The parameter \(C\) therefore sets the balance between the two goals. With a
large \(C\), mistakes are expensive and the problem behaves like the
hard-margin SVM. With a small \(C\), a large margin matters more, and more
mistakes are tolerated.

\begin{moral}
    \(\hlr{C}\): trade-off of \green{large margin} \& \hlb{noise
    tolerance}.
\end{moral}

\subsection{Soft-Margin SVM (1/2)}

The two kinds of constraint can be written as a single line. When example
\(n\) is classified correctly, the indicator is \(0\) and the right-hand side
is \(1\). When it is wrong, the indicator is \(1\) and the right-hand side
is \(1 - \infty = -\infty\):
\[
    \begin{aligned}
        \min_{b, \bm{w}} \quad & \green{\frac{1}{2} \bm{w}^\top \bm{w}} +
            \hlr{C} \cdot \sum_{n=1}^N \hlb{\llbracket y_n \neq
            \sign(\bm{w}^\top \bm{z}_n + b) \rrbracket}\\
        \text{s.t.} \quad & y_n (\bm{w}^\top \bm{z}_n + b) \ge 1 - \infty
            \cdot \hlb{\llbracket y_n \neq \sign(\bm{w}^\top \bm{z}_n + b)
            \rrbracket}
    \end{aligned}
\]
This says exactly what we want, but it has two problems.
\begin{itemize}
    \item \(\hlb{\llbracket \cdot \rrbracket}\) is non-linear, so the
        problem is \hlo{\textbf{not QP anymore :-(}} --- what about the dual?
        the kernel?
    \item It \textbf{cannot distinguish small error} (slightly away from the
        fat boundary) \textbf{or large error} (a\ldots w\ldots a\ldots
        y\ldots\ from the fat boundary).
\end{itemize}

The first problem is serious. A QP needs a quadratic objective and linear
constraints. The indicator jumps between \(0\) and \(1\), so it is neither,
and a QP solver cannot handle it. The dual and the kernel trick of the last
two chapters were both built on top of a QP, so we would lose them too.

The second problem is about fairness. The indicator only asks `right or
wrong?'. An example that is barely on the wrong side costs the same \(C\) as
one that lies deep inside the other class.

One change fixes both problems: instead of asking \emph{whether} an example
is wrong, measure \emph{how far} it falls short of the margin.
\begin{itemize}
    \item Record the `\vocab{margin violation}' of example \(n\) by a new
        variable \(\hlp{\xi_n}\): require \(y_n (\bm{w}^\top \bm{z}_n + b)
        \ge 1 - \hlp{\xi_n}\) with \(\hlp{\xi_n} \ge 0\). These are
        \hlo{\textbf{linear constraints}}.
    \item Penalize with \hlp{\textbf{margin violation}} instead of
        \hlb{\textbf{error count}}. The penalty \(\hlr{C} \cdot \sum_n
        \hlp{\xi_n}\) is linear, so we still have a \hlo{\textbf{quadratic
        objective}}.
\end{itemize}

\begin{definition}[Soft-margin SVM]
    The \vocab{soft-margin SVM} solves
    \[
        \begin{aligned}
            \min_{b, \bm{w}, \hlp{\bm{\xi}}} \quad & \green{\frac{1}{2}
                \bm{w}^\top \bm{w}} + \hlr{C} \cdot \sum_{n=1}^N \hlp{\xi_n}\\
            \text{s.t.} \quad & y_n (\bm{w}^\top \bm{z}_n + b) \ge 1 -
                \hlp{\xi_n} \text{ and } \hlp{\xi_n} \ge 0 \text{ for all }
                n.
        \end{aligned}
    \]
\end{definition}

Think of \(\xi_n\) as the amount by which example \(n\) is allowed to fall
short of the margin.
\begin{itemize}
    \item If the example already satisfies \(y_n (\bm{w}^\top \bm{z}_n + b)
        \ge 1\), we can take \(\xi_n = 0\), and it costs nothing.
    \item If it does not, \(\xi_n\) has to be at least \(1 - y_n
        (\bm{w}^\top \bm{z}_n + b)\). The further the example is from where
        it should be, the larger \(\xi_n\), and the larger the penalty \(C
        \xi_n\).
\end{itemize}
Since we minimize, the optimal \(\xi_n\) is always the smallest value
allowed: \(\xi_n = \max\bigl(0, 1 - y_n (\bm{w}^\top \bm{z}_n +
b)\bigr)\). Small violations are now cheap and large ones are expensive.

\subsection{Soft-Margin SVM (2/2)}\label{sec:soft-geometry}

\begin{center}
    \begin{minipage}[c]{0.56\textwidth}
        \begin{itemize}
            \item record `\hlp{\textbf{margin violation}}' by
                \(\hlp{\xi_n}\)
            \item penalize with \hlp{\textbf{margin violation}}
        \end{itemize}
        \[
            \begin{aligned}
                \min_{b, \bm{w}, \hlp{\bm{\xi}}} \quad & \green{\frac{1}{2}
                    \bm{w}^\top \bm{w}} + \hlr{C} \cdot \sum_{n=1}^N
                    \hlp{\xi_n}\\
                \text{s.t.} \quad & y_n (\bm{w}^\top \bm{z}_n + b) \ge 1 -
                    \hlp{\xi_n}\\
                & \text{and } \hlp{\xi_n} \ge 0 \text{ for all } n
            \end{aligned}
        \]
    \end{minipage}%
    \hfill
    \begin{minipage}[c]{0.4\textwidth}
        \centering
        \begin{tikzpicture}[scale=4]
            \begin{scope}
                \clip (0,0) rectangle (1,1);
                % centre line through (0.08,0) with direction (0.81,1);
                % the fat band has half-width 0.144 around it
                \fill[fatband] (-0.267,-0.2) -- (0.867,1.2)
                    -- (1.237,1.2) -- (0.103,-0.2) -- cycle;
                \draw[sepline] (0.08,0) -- (0.89,1);
            \end{scope}
            \draw[plotbox] (0,0) rectangle (1,1);
            \foreach \p in {(0.141,0.908), (0.408,0.919), (0.528,0.877),
                (0.373,0.754), (0.106,0.585), (0.123,0.349)}
                {\node[dneg] at \p {\(\times\)};}
            \foreach \p in {(0.887,0.585), (0.496,0.440), (0.496,0.254),
                (0.641,0.113), (0.873,0.127)}{\node[dpos] at \p {};}
            \draw[-{Stealth[length=1.3mm]}, violet!80!black, thick]
                (0.505,0.433) -- (0.571,0.379);
            \node[font=\scriptsize\bfseries, text=violet!80!black, anchor=west]
                at (0.585,0.395) {violation};
        \end{tikzpicture}
    \end{minipage}
\end{center}

In the picture the grey strip is the fat boundary, and the circle inside it
has a violation \(\xi_n > 0\). Before explaining the strip, we check what
the \(1\) in the constraint means now.

\paragraph{The \(1\) is still a choice of scale.} In
\autoref{sec:standard-lm}, \((b, \bm{w})\) and \(c \cdot (b, \bm{w})\) with
\(c > 0\) give the same hyperplane, so we were free to fix the scale by
\(\min_n y_n (\bm{w}^\top \bm{x}_n + b) = 1\) (\autoref{lem:special-scaling}).
The soft-margin constraint keeps this \(1\), but it can no longer mean
\(\min_n y_n (\bm{w}^\top \bm{z}_n + b) = 1\): an example with \(\xi_n > 0\)
has \(y_n (\bm{w}^\top \bm{z}_n + b) < 1\). Instead, the \(1\) fixes where
the two edges of the strip are: at \(\bm{w}^\top \bm{z} + b = \pm 1\).

Does this rule out any candidate? No. Take any hyperplane and any strip of
half-width \(h\) around it. Write the hyperplane as \(\bm{w}^\top \bm{z} + b
= 0\) and scale \((b, \bm{w})\) so that \(\norm{\bm{w}} = 1/h\). Then the
edges \(\bm{w}^\top \bm{z} + b = \pm 1\) are exactly at distance \(h\). So
every pair (hyperplane, strip width) is some \((b, \bm{w})\), and the primal
searches over all of them.

What changes is the role of the scale. In the hard-margin SVM, \((b,
\bm{w})\) and \(c \cdot (b, \bm{w})\) were the same candidate. Here they
have the same hyperplane but different strips (half-width \(1 /
\norm{\bm{w}}\) versus \(1 / (c \norm{\bm{w}})\)), hence different
\(\xi_n\) and a different objective value. The scale of \((b, \bm{w})\) is
no longer free: it is how the primal chooses the width of the strip.

Why is the strip there, and why does the size of \(\xi_n\) tell us where an
example lies? We answer this in three steps.

\paragraph{Step 1: the strip comes from the hyperplane.} Any \((b,
\bm{w})\) gives a hyperplane \(\bm{w}^\top \bm{z} + b = 0\) and two parallel
hyperplanes \(\bm{w}^\top \bm{z} + b = +1\) and \(\bm{w}^\top \bm{z} + b =
-1\), one on each side. The fat boundary is the strip between these two,
where \(\abs{\bm{w}^\top \bm{z} + b} < 1\). This is pure geometry: it holds
for any \((b, \bm{w})\), whichever problem produced it. For the soft-margin
SVM we use its optimal \((b, \bm{w})\).

\paragraph{Step 2: \(y_n (\bm{w}^\top \bm{z}_n + b)\) says where
\(\bm{z}_n\) is.} By \autoref{sec:standard-lm}, the distance from
\(\bm{z}_n\) to the hyperplane is \(\abs{\bm{w}^\top \bm{z}_n + b} /
\norm{\bm{w}}\). Multiplying by \(y_n\) adds a sign: \(y_n (\bm{w}^\top
\bm{z}_n + b)\) is positive when \(\bm{z}_n\) is on the side of its own
label, and negative when it is on the other side. So \(y_n (\bm{w}^\top
\bm{z}_n + b) / \norm{\bm{w}}\) is a \emph{signed} distance, and the edge of
the strip on the example's own side is where this number equals \(1 /
\norm{\bm{w}}\), i.e.\ where \(y_n (\bm{w}^\top \bm{z}_n + b) = 1\). This
gives:
\begin{itemize}
    \item \(y_n (\bm{w}^\top \bm{z}_n + b) \ge 1\): on or beyond its own
        edge, i.e.\ outside the strip (or on its edge), correct side;
    \item \(0 < y_n (\bm{w}^\top \bm{z}_n + b) < 1\): between the hyperplane
        and its own edge, i.e.\ inside the strip, correct side;
    \item \(y_n (\bm{w}^\top \bm{z}_n + b) = 0\): on the hyperplane;
    \item \(y_n (\bm{w}^\top \bm{z}_n + b) < 0\): on the other side, i.e.\
        misclassified.
\end{itemize}

\paragraph{Step 3: \(\xi_n\) turns this into a violation.} We saw that at
the optimum \(\xi_n = \max\bigl(0, 1 - y_n (\bm{w}^\top \bm{z}_n +
b)\bigr)\). So \(\xi_n\) is how far \(y_n (\bm{w}^\top \bm{z}_n + b)\) falls
below \(1\), and each range of \(\xi_n\) matches one line of Step 2:
\begin{itemize}
    \item \(\xi_n = 0\) \(\iff\) \(y_n (\bm{w}^\top \bm{z}_n + b) \ge 1\):
        outside the strip or on its edge;
    \item \(0 < \xi_n < 1\) \(\iff\) \(0 < y_n (\bm{w}^\top \bm{z}_n + b) <
        1\): inside the strip, still correct (the circle in the picture);
    \item \(\xi_n = 1\) \(\iff\) \(y_n (\bm{w}^\top \bm{z}_n + b) = 0\): on
        the hyperplane;
    \item \(\xi_n > 1\) \(\iff\) \(y_n (\bm{w}^\top \bm{z}_n + b) < 0\):
        misclassified.
\end{itemize}
Measured as a distance, the gap between the example and its own edge is
\(\xi_n / \norm{\bm{w}}\); this is the arrow in the picture.

How does this differ from the hard-margin SVM? In the hard-margin SVM,
\(y_n (\bm{w}^\top \bm{z}_n + b) \ge 1\) is required for every example. So
no example lies inside the strip, and its edges pass through the closest
examples. Here examples may lie inside
the strip. The edges are no longer `where the closest examples are', but
the lines beyond which an example pays nothing. The half-width \(1 /
\norm{\bm{w}}\) of the strip is chosen by the SVM itself, which trades a
wide strip against the total violation.

\begin{itemize}
    \item The parameter \(\hlr{C}\) is the trade-off of \green{large margin}
        \& \hlp{margin violation}:
        \begin{itemize}
            \item large \(\hlr{C}\): want less \hlp{margin violation};
            \item small \(\hlr{C}\): want \green{large margin}.
        \end{itemize}
    \item The problem is a \hlo{QP} of \(\hlo{\tilde{d}} + 1 + N\)
        variables, \(2N\) constraints.
\end{itemize}

When \(C\) is very large, any violation becomes too expensive, and we are
back to the hard-margin SVM. To count the size of the QP: the variables are
\(b\), the \(\tilde{d}\) entries of \(\bm{w}\), and one \(\xi_n\) for each
example; the constraints are the \(N\) margin constraints plus the \(N\)
constraints \(\xi_n \ge 0\).

The number of variables still contains \(\tilde{d}\). In
\autoref{sec:lagrange-dual} we removed \(\tilde{d}\) by moving to the dual
problem, so the natural next step is to try the same here.

\begin{question}
    Next: remove dependence on \(\hlo{\tilde{d}}\) by soft-margin SVM primal
    \(\Rightarrow\) \hlo{\textbf{dual}}?
\end{question}

\Needspace{8\baselineskip}
\begin{exercise}[Fun Time]
    At the optimal solution of
    \[
        \begin{aligned}
            \min_{b, \bm{w}, \bm{\xi}} \quad & \frac{1}{2} \bm{w}^\top
                \bm{w} + C \cdot \sum_{n=1}^N \xi_n\\
            \text{s.t.} \quad & y_n (\bm{w}^\top \bm{z}_n + b) \ge 1 - \xi_n
                \text{ and } \xi_n \ge 0 \text{ for all } n,
        \end{aligned}
    \]
    assume that \(y_1 (\bm{w}^\top \bm{z}_1 + b) = -10\). What is the
    corresponding \(\xi_1\)?
    \begin{enumerate}[(1)]
        \item \(1\);
        \item \(11\);
        \item \(21\);
        \item \(31\).
    \end{enumerate}
\end{exercise}

\begin{proof}[Reference Answer]
    (2). \(\xi_1\) is simply \(1 - y_1 (\bm{w}^\top \bm{z}_1 + b)\) when
    \(y_1 (\bm{w}^\top \bm{z}_1 + b) \le 1\). Here that is \(1 - (-10) =
    11\). This is the formula \(\xi_n = \max\bigl(0, 1 - y_n (\bm{w}^\top
    \bm{z}_n + b)\bigr)\) from above.
\end{proof}

\section{Dual Problem}\label{sec:soft-dual}

\source{slides p.7--11}

We now repeat the steps of \autoref{sec:lagrange-dual} for the soft-margin
problem. Almost everything goes through as before. The only new things are
the variables \(\xi_n\) and their constraints \(\xi_n \ge 0\).

\subsection{Lagrange Dual}

\[
    \begin{aligned}
        \text{primal:} \quad \min_{b, \bm{w}, \bm{\xi}} \quad & \frac{1}{2}
            \bm{w}^\top \bm{w} + C \cdot \sum_{n=1}^N \xi_n\\
        \text{s.t.} \quad & \dvar{y_n (\bm{w}^\top \bm{z}_n + b) \ge 1 -
            \xi_n} \text{ and } \hlb{\xi_n \ge 0} \text{ for all } n
    \end{aligned}
\]

The primal has two kinds of constraint, so we need two kinds of Lagrange
multiplier:
\begin{itemize}
    \item \(\dvar{\alpha_n}\) for the margin constraint (red), as in the
        hard-margin case;
    \item \(\hlb{\beta_n}\) for the new constraint \(\hlb{\xi_n \ge 0}\)
        (blue).
\end{itemize}
As in \autoref{sec:lagrange-dual}, we write each constraint in the form
`something \(\le 0\)', multiply it by its multiplier, and add it to the
objective. The margin constraint becomes \(1 - \xi_n - y_n (\bm{w}^\top
\bm{z}_n + b) \le 0\), and \(\xi_n \ge 0\) becomes \(-\xi_n \le 0\). This
gives the Lagrange function with Lagrange multipliers \(\dvar{\alpha_n}\) and
\(\hlb{\beta_n}\):
\[
    \begin{aligned}
        \mathcal{L}(b, \bm{w}, \bm{\xi}, \dvar{\bm{\alpha}},
            \hlb{\bm{\beta}}) &= \frac{1}{2} \bm{w}^\top \bm{w} + C \cdot
            \sum_{n=1}^N \xi_n\\
        &\quad \dvar{{} + \sum_{n=1}^N \alpha_n \cdot \bigl( 1 - \xi_n - y_n
            (\bm{w}^\top \bm{z}_n + b) \bigr)} \hlb{{} + \sum_{n=1}^N
            \beta_n \cdot (-\xi_n)}.
    \end{aligned}
\]

The reasoning is the same as in \autoref{lem:svm-minmax}. If we maximize
\(\mathcal{L}\) over all \(\alpha_n \ge 0\) and \(\beta_n \ge 0\), the
result is \(+\infty\) whenever some constraint is broken, and it is just the
primal objective when all constraints hold. So minimizing that maximum over
\((b, \bm{w}, \bm{\xi})\) is the primal problem. The primal is a convex QP,
and it always has a feasible point (take any \((b, \bm{w})\) and make every
\(\xi_n\) large enough), so by strong duality (\autoref{fact:strong-duality})
we may swap \(\min\) and \(\max\). What we want is the Lagrange dual
\[
    \max_{\dvar{\alpha_n \ge 0}, \; \hlb{\beta_n \ge 0}} \left( \min_{b,
    \bm{w}, \bm{\xi}} \mathcal{L}(b, \bm{w}, \bm{\xi}, \dvar{\bm{\alpha}},
    \hlb{\bm{\beta}}) \right).
\]

\subsection{Simplify \texorpdfstring{\(\xi_n\)}{xi} and \texorpdfstring{\(\hlb{\beta_n}\)}{beta}}

\[
    \max_{\dvar{\alpha_n \ge 0}, \; \hlb{\beta_n \ge 0}} \left( \min_{b,
    \bm{w}, \bm{\xi}} \frac{1}{2} \bm{w}^\top \bm{w} + C \cdot \sum_{n=1}^N
    \xi_n \dvar{{} + \sum_{n=1}^N \alpha_n \cdot \bigl( 1 - \xi_n - y_n
    (\bm{w}^\top \bm{z}_n + b) \bigr)} \hlb{{} + \sum_{n=1}^N \beta_n \cdot
    (-\xi_n)} \right)
\]

The inner \(\min\) is over \(b\), \(\bm{w}\) and \(\bm{\xi}\), with no
constraints. We deal with \(\xi_n\) first, in the same way that
\autoref{sec:lagrange-dual} dealt with \(b\).
\begin{itemize}
    \item The inner problem is unconstrained, so at its optimum
        \(\frac{\partial \mathcal{L}}{\partial \xi_n} = 0 = C -
        \dvar{\alpha_n} - \hlb{\beta_n}\).
    \item There is no loss of optimality if solving with the
        \textbf{implicit} constraint \(\hlb{\beta_n} = C - \dvar{\alpha_n}\)
        and the \textbf{explicit} constraint \(0 \le \dvar{\alpha_n} \le
        \hlb{C}\): \(\hlb{\beta_n}\) is removed.
\end{itemize}

Why may we add this constraint? Collect the terms of \(\mathcal{L}\) that
contain \(\xi_n\): they are \(C \xi_n - \alpha_n \xi_n - \beta_n \xi_n = (C
- \alpha_n - \beta_n) \xi_n\). This is a straight line in \(\xi_n\).
\begin{itemize}
    \item If its slope \(C - \alpha_n - \beta_n\) is not zero, we can make
        \(\mathcal{L}\) as small as we like by moving \(\xi_n\) to \(+\infty\)
        or \(-\infty\). The inner \(\min\) is then \(-\infty\), and the outer
        \(\max\) would never choose such \((\bm{\alpha}, \bm{\beta})\).
    \item So we lose nothing by keeping only the multipliers with \(C -
        \alpha_n - \beta_n = 0\), i.e.\ \(\beta_n = C - \alpha_n\).
\end{itemize}
Once \(\beta_n = C - \alpha_n\), the condition \(\beta_n \ge 0\) becomes
\(\alpha_n \le C\). Together with \(\alpha_n \ge 0\) this gives \(0 \le
\alpha_n \le C\). After that, \(\beta_n\) no longer carries any information
of its own, so it disappears from the problem.

\hlp{\(\bm{\xi}\) \textbf{can also be removed :-)}}, \hlp{like how we
removed \(b\)}. With \(C - \alpha_n - \beta_n = 0\), every term that
contains \(\xi_n\) has coefficient zero, so the whole last sum drops out:
\[
    \max_{0 \le \dvar{\alpha_n} \le C, \; \hlb{\beta_n} = C -
    \dvar{\alpha_n}} \left( \min_{b, \bm{w}, \bm{\xi}} \frac{1}{2}
    \bm{w}^\top \bm{w} + \sum_{n=1}^N \dvar{\alpha_n \bigl( 1 - y_n
    (\bm{w}^\top \bm{z}_n + b) \bigr)} + \cancel{\sum_{n=1}^N (C -
    \dvar{\alpha_n} - \hlb{\beta_n}) \cdot \xi_n} \right)
\]

\subsection{Other Simplifications}

\[
    \max_{0 \le \dvar{\alpha_n} \le C, \; \hlb{\beta_n} = C -
    \dvar{\alpha_n}} \left( \min_{b, \bm{w}} \frac{1}{2} \bm{w}^\top \bm{w}
    + \sum_{n=1}^N \dvar{\alpha_n \bigl( 1 - y_n (\bm{w}^\top \bm{z}_n + b)
    \bigr)} \right)
\]

\hlp{\textbf{Familiar? :-)}}
\begin{itemize}
    \item The inner problem is the \hlp{\textbf{same as hard-margin SVM}}
        (\autoref{sec:lagrange-dual}).
    \item \(\frac{\partial \mathcal{L}}{\partial b} = 0\): no loss of
        optimality if solving with constraint \(\sum_{n=1}^N
        \dvar{\alpha_n} y_n = 0\).
    \item \(\frac{\partial \mathcal{L}}{\partial w_i} = 0\): no loss of
        optimality if solving with constraint \(\bm{w} = \sum_{n=1}^N
        \dvar{\alpha_n} y_n \bm{z}_n\).
\end{itemize}

Compare this with \autoref{sec:lagrange-dual}. The expression inside the
\(\min\) is exactly the hard-margin one. The only difference is outside: each
\(\alpha_n\) now has the upper limit \(C\). Therefore the standard dual can
be derived \hlp{using the same steps as \autoref{sec:lagrange-dual}}. Using
the two constraints above, \(b\) drops out, and \(\sum_n \alpha_n y_n
\bm{w}^\top \bm{z}_n = \bm{w}^\top \bm{w}\). The objective becomes
\(-\frac{1}{2} \norm{\sum_n \alpha_n y_n \bm{z}_n}^2 + \sum_n \alpha_n\).
Multiplying by \(-1\) turns the \(\max\) into a \(\min\).

\subsection{Standard Soft-Margin SVM Dual}

\begin{definition}[Standard soft-margin SVM dual]
    \[
        \begin{aligned}
            \min_{\dvar{\bm{\alpha}}} \quad & \frac{1}{2} \sum_{n=1}^N
                \sum_{m=1}^N \dvar{\alpha_n \alpha_m} y_n y_m \bm{z}_n^\top
                \bm{z}_m - \sum_{n=1}^N \dvar{\alpha_n}\\
            \text{subject to} \quad & \sum_{n=1}^N y_n \dvar{\alpha_n} = 0;\\
            & 0 \le \dvar{\alpha_n} \le C, \text{ for } n = 1, 2, \dots ,
                N;\\
            \text{implicitly} \quad & \bm{w} = \sum_{n=1}^N \dvar{\alpha_n}
                y_n \bm{z}_n;\\
            & \hlb{\beta_n} = C - \dvar{\alpha_n}, \text{ for } n = 1, 2,
                \dots , N.
        \end{aligned}
    \]
\end{definition}

The only difference to hard-margin is the \hlb{upper bound} on
\(\dvar{\alpha_n}\). This upper bound is what makes the problem `soft'. In
the hard-margin dual, an example that cannot be separated would push its
\(\alpha_n\) towards infinity, and the dual would have no finite solution.
Here no \(\alpha_n\) can exceed \(C\), so no single example can take over
the solution.

\begin{moral}
    Another (convex) \hlo{\textbf{QP}}, with \(\hlo{N}\) variables \&
    \(\hlo{2N + 1}\) constraints.
\end{moral}

The \(2N + 1\) constraints are one equality \(\sum_n y_n \alpha_n = 0\),
\(N\) lower bounds \(\alpha_n \ge 0\), and \(N\) upper bounds \(\alpha_n \le
C\). The number of variables no longer depends on \(\tilde{d}\), which was
the goal.

\begin{exercise}[Fun Time]
    In the soft-margin SVM, assume that we want to increase the parameter
    \(C\) by \(2\). How shall the corresponding dual problem be changed?
    \begin{enumerate}[(1)]
        \item the upper bound of \(\alpha_n\) shall be halved;
        \item the upper bound of \(\alpha_n\) shall be decreased by \(2\);
        \item the upper bound of \(\alpha_n\) shall be increased by \(2\);
        \item the upper bound of \(\alpha_n\) shall be doubled.
    \end{enumerate}
\end{exercise}

\begin{proof}[Reference Answer]
    (3). Because \(C\) is exactly the upper bound of \(\alpha_n\), increasing
    \(C\) by \(2\) in the primal problem is equivalent to increasing the
    upper bound by \(2\) in the dual problem.
\end{proof}

\section{Messages behind Soft-Margin SVM}\label{sec:soft-messages}

\source{slides p.12--16}

We have the dual. This section turns it into an algorithm, works out how to
compute \(b\), and explains what the solution \(\alpha_n\) tells us about
each example.

\subsection{Kernel Soft-Margin SVM}

In the dual, \(\bm{z}_n\) appears only in inner products
\(\bm{z}_n^\top \bm{z}_m\). So, exactly as in \autoref{sec:kernel-trick}, we
can replace each of them by a kernel \(K(\bm{x}_n, \bm{x}_m)\).

\begin{algorithm}[H]
    \caption{\hlp{\textbf{Kernel}} Soft-Margin \hlo{SV}M}\label{alg:kernel-soft-svm}
    \(\dvar{q_{n,m}} = y_n y_m \hlp{K}(\bm{x}_n, \bm{x}_m)\);
        \(\dvar{\bm{p}} = -\bm{1}_N\); \(\green{(\mathrm{A}, \bm{c})}\) for
        \green{equ./lower-bound/}\hlb{\textbf{upper-bound}}
        \green{constraints}\;
    \(\bm{\alpha} \gets \mathrm{QP}(\dvar{\mathrm{Q}_{\mathrm{D}}},
        \dvar{\bm{p}}, \green{\mathrm{A}}, \green{\bm{c}})\)\;
    \(b \gets\) ?\;
    \Return\ \hlo{SVs} and their \(\alpha_n\) as well as \(b\) such that for
        new \(\bm{x}\),
        \(g_{\mathrm{SVM}}(\bm{x}) = \sign \left( \sum_{\hlo{\text{SV}}
        \text{ indices } n} \alpha_n \hlo{y_n} \hlp{K}(\hlo{\bm{x}_n},
        \bm{x}) + b \right)\)\;
\end{algorithm}

\begin{itemize}
    \item It is \hlb{almost} the same as hard-margin
        (\autoref{alg:kernel-hard-svm}).
    \item It is more flexible than hard-margin: the \textbf{primal/dual}
        are \textbf{always solvable}.
\end{itemize}

The only change from \autoref{alg:kernel-hard-svm} is in line~1: we add
\(N\) more constraints for the upper bounds \(\alpha_n \le C\). The second
point is a real gain. If the data cannot be separated in \(\mathcal{Z}\),
the hard-margin problem has no solution at all. The soft-margin problem
always has one: the primal always has a feasible point, as we saw, and in
the dual every \(\alpha_n\) is confined to the interval \([0, C]\).

\begin{question}
    Remaining question: \hlo{step \textcircled{3}}?
\end{question}

\subsection{Solving for \texorpdfstring{\(b\)}{b}}\label{sec:soft-solve-b}

In the hard-margin case we got \(b\) from complementary slackness
(\autoref{sec:solving-dual}). The soft-margin problem has one such condition
for each kind of multiplier. The first is for \(\alpha_n\). The second
comes from \(\beta_n \xi_n = 0\); since \(\beta_n = C - \alpha_n\), it reads
\((C - \alpha_n) \xi_n = 0\).

\begin{center}
    \begin{tabular}{@{}l@{\qquad}l@{}}
        \textbf{hard-margin SVM} & \textbf{soft-margin SVM}\\[2pt]
        complementary slackness: & complementary slackness:\\
        \(\dvar{\alpha_n} \bigl( 1 - y_n (\bm{w}^\top \bm{z}_n + b) \bigr)
            = 0\) &
            \(\dvar{\alpha_n} \bigl( 1 - \hlp{\xi_n} - y_n (\bm{w}^\top
            \bm{z}_n + b) \bigr) = 0\)\\
        & \((C - \dvar{\alpha_n}) \hlp{\xi_n} = 0\)\\[4pt]
        \textbullet\ SV (\(\dvar{\alpha_s} > 0\)) & \textbullet\ SV
            (\(\dvar{\alpha_s} > 0\))\\
        \quad \(\Rightarrow b = y_s - \bm{w}^\top \bm{z}_s\) & \quad
            \(\Rightarrow b = y_s - y_s \hlp{\xi_s} - \bm{w}^\top
            \bm{z}_s\)\\
        & \textbullet\ \textbf{free} (\(\dvar{\alpha_s} < C\))\\
        & \quad \(\Rightarrow \hlp{\xi_s} = 0\)
    \end{tabular}
\end{center}

Let us go through the right-hand column step by step.
\begin{enumerate}
    \item Take an SV, i.e.\ an example \(s\) with \(\alpha_s > 0\). The
        first condition then forces the bracket to be zero: \(y_s
        (\bm{w}^\top \bm{z}_s + b) = 1 - \xi_s\).
    \item Multiply both sides by \(y_s\) (and use \(y_s^2 = 1\)). This
        gives \(b = y_s - y_s \xi_s - \bm{w}^\top \bm{z}_s\).
    \item This formula contains \(\xi_s\), which we do not know. The second
        condition helps here: if also \(\alpha_s < C\), then \(C - \alpha_s
        > 0\), so \(\xi_s\) must be \(0\).
\end{enumerate}
An SV with \(0 < \alpha_s < C\) is called a \vocab{free} SV. For a free SV
the formula becomes \(b = y_s - \bm{w}^\top \bm{z}_s\), and we compute
\(\bm{w}^\top \bm{z}_s\) with the kernel.

\begin{moral}
    Solve unique \(b\) with \hlp{free SV \((\bm{x}_s, y_s)\)}:
    \[
        b = \hlp{y_s} - \sum_{\hlo{\text{SV}} \text{ indices } n}
        \dvar{\alpha_n} y_n K(\bm{x}_n, \hlp{\bm{x}_s})
    \]
    \gone{--- range of \(b\) otherwise.}
\end{moral}

What does `range of \(b\) otherwise' mean? If there is no free SV, every
\(\alpha_n\) is either \(0\) or \(C\), and the conditions above give only
inequalities for \(b\), not an equation:
\begin{itemize}
    \item an example with \(\alpha_n = 0\) has \(\xi_n = 0\), so \(y_n
        (\bm{w}^\top \bm{z}_n + b) \ge 1\);
    \item an example with \(\alpha_n = C\) has \(y_n (\bm{w}^\top \bm{z}_n +
        b) = 1 - \xi_n \le 1\).
\end{itemize}
Each inequality is a lower or an upper bound on \(b\). Together they give an
interval.

Why is every \(b\) in this interval optimal? Keep the optimal
\(\bm{\alpha}\), and with it \(\bm{w}\). Pick any \(b\) in the interval and
set \(\xi_n = \max\bigl(0, 1 - y_n (\bm{w}^\top \bm{z}_n + b)\bigr)\). Then
all the KKT conditions hold.
\begin{itemize}
    \item The primal constraints hold because of how \(\xi_n\) was chosen.
    \item The conditions that do not involve \(b\) --- \(0 \le \alpha_n \le
        C\), \(\sum_n y_n \alpha_n = 0\), \(\bm{w} = \sum_n \alpha_n y_n
        \bm{z}_n\) and \(\beta_n = C - \alpha_n\) --- are already met by the
        optimal \(\bm{\alpha}\).
    \item The two complementary slackness conditions are exactly the
        inequalities above. If \(\alpha_n = 0\), the bound \(y_n
        (\bm{w}^\top \bm{z}_n + b) \ge 1\) makes \(\xi_n = 0\). If \(\alpha_n
        = C\), the bound \(y_n (\bm{w}^\top \bm{z}_n + b) \le 1\) makes
        \(\xi_n = 1 - y_n (\bm{w}^\top \bm{z}_n + b)\), so the bracket in
        the first condition is zero.
\end{itemize}
For a convex QP like this one, the KKT conditions are not only necessary
but also sufficient; the proof after \autoref{fact:kkt} goes through
unchanged. So \((b, \bm{w}, \bm{\xi})\) is optimal. A \(b\) outside the
interval breaks one of the inequalities, so it is not.

This case is rare in practice; SVM software then simply picks a value inside
the interval.


\subsection{Soft-Margin Gaussian SVM in Action}

The three pictures show a soft-margin Gaussian SVM with the same \(\gamma\)
and three different values of \(C\). In each picture, blue is the region
predicted \(+1\) and red the region predicted \(-1\). The darker strip around
the boundary is the fat boundary; any example inside it has \(\xi_n > 0\).

\begin{note}
    The data are the 40 examples of the Gaussian SVM pictures in
    \autoref{sec:gaussian-kernel}, with three labels flipped to act as noise:
    two circles became crosses (one in the middle of the circles, one on the
    right edge), and one cross on the lower right became a circle. We
    measured the points from the slides and solved the SVMs ourselves.
\end{note}

\begin{center}
    \begin{tabular}{@{}c@{\hspace{1.2em}}c@{\hspace{1.2em}}c@{}}
        \ksvmpanel[3.6]{\smsvmBi}{\smsvmData}{} &
        \ksvmpanel[3.6]{\smsvmBii}{\smsvmData}{} &
        \ksvmpanel[3.6]{\smsvmBiii}{\smsvmData}{}\\[4pt]
        \(C = 1\) & \(C = 10\) & \(C = 100\)
    \end{tabular}
\end{center}

Read the pictures from left to right.
\begin{itemize}
    \item \(C = 1\): violations are cheap. The boundary is smooth and the
        strip is wide. Three training examples are misclassified, including
        the lone cross among the circles and the circle on the lower right.
    \item \(C = 10\): violations cost more, so the boundary bends to fix
        some of them. One training error is left.
    \item \(C = 100\): there are no training errors. The boundary even draws
        a small island around the lone cross among the circles.
\end{itemize}
The cross inside that island is one of the noisy labels, so the right-hand
picture is fitting noise.
\begin{itemize}
    \item Large \(C\) \(\Longrightarrow\) less \hlb{noise tolerance}
        \(\Longrightarrow\) \hlo{`overfit'?}
    \item \hlo{\textbf{Warning: SVM can still overfit :-(}}
\end{itemize}

In other words, \(C\) behaves like \(\gamma\) in
\autoref{sec:gaussian-kernel}: if it is too large, the SVM overfits. Both
parameters have to be chosen with care.

\begin{moral}
    Soft-margin Gaussian SVM: need \hlo{careful selection of \((\gamma,
    C)\)}.
\end{moral}

\subsection{Physical Meaning of \texorpdfstring{\(\alpha_n\)}{alpha}}

Where an example lies relative to the fat boundary is decided by the number
\(y_n (\bm{w}^\top \bm{z}_n + b)\), as explained in
\autoref{sec:soft-geometry}: \(\ge 1\) means outside the strip or on its
edge, \(= 1\) exactly on the edge, and \(< 1\) inside the strip or on the
wrong side.
We could compute this number for every example. But we do not have to: the
two complementary slackness conditions
\[
    \dvar{\alpha_n} \hlo{\bigl( 1 - \hlp{\xi_n} - y_n (\bm{w}^\top \bm{z}_n +
    b) \bigr)} = 0, \qquad (C - \dvar{\alpha_n}) \hlp{\xi_n} = 0
\]
connect it to \(\alpha_n\), which the dual has already given us. So the
value of \(\alpha_n\) alone tells us which of these cases are possible.

\begin{center}
    \begin{minipage}[c]{0.56\textwidth}
        \begin{itemize}
            \item \textbf{non SV} (\(0 = \dvar{\alpha_n}\)):
                \(\hlp{\xi_n} = 0\), `away from'/\gone{on} \hlo{fat
                boundary};
            \item \(\square\) \textbf{free SV} (\(0 < \dvar{\alpha_n} <
                C\)): \(\hlp{\xi_n} = 0\), on \hlo{fat boundary}, locates
                \(b\);
            \item \(\triangle\) \textbf{bounded SV} (\(\dvar{\alpha_n} =
                C\)): \(\hlp{\xi_n} =\) violation amount,
                `violate'/\gone{on} \hlo{fat boundary}.
        \end{itemize}
    \end{minipage}%
    \hfill
    \begin{minipage}[c]{0.4\textwidth}
        \centering
        \softAlphaPanel
    \end{minipage}
\end{center}

Why does each line hold? In every case we use the second condition to find
\(\xi_n\), and then the first condition (or the primal constraint \(y_n
(\bm{w}^\top \bm{z}_n + b) \ge 1 - \xi_n\)) to find \(y_n (\bm{w}^\top
\bm{z}_n + b)\).
\begin{itemize}
    \item \textbf{Non SV, \(\alpha_n = 0\).}
        \begin{itemize}
            \item Second condition: \(C - \alpha_n = C > 0\), so \(\xi_n =
                0\).
            \item First condition: it reads \(0 \cdot (\ldots) = 0\), which
                says nothing. We use the primal constraint instead: \(y_n
                (\bm{w}^\top \bm{z}_n + b) \ge 1 - \xi_n = 1\).
            \item So the example is outside the strip, or exactly on its
                edge.
        \end{itemize}
    \item \textbf{Free SV, \(0 < \alpha_n < C\).}
        \begin{itemize}
            \item Second condition: \(C - \alpha_n > 0\), so \(\xi_n = 0\).
            \item First condition: \(\alpha_n > 0\), so the bracket must be
                zero: \(y_n (\bm{w}^\top \bm{z}_n + b) = 1 - \xi_n = 1\).
            \item So the example is exactly on the edge of the strip. This
                is an equation, not an inequality, which is why a free SV
                pins down \(b\) (\autoref{sec:soft-solve-b}).
        \end{itemize}
    \item \textbf{Bounded SV, \(\alpha_n = C\).}
        \begin{itemize}
            \item Second condition: \(C - \alpha_n = 0\), so it reads \(0
                \cdot \xi_n = 0\) and says nothing about \(\xi_n\).
            \item First condition: \(\alpha_n = C > 0\), so the bracket must
                be zero: \(\xi_n = 1 - y_n (\bm{w}^\top \bm{z}_n + b)\).
            \item Since \(\xi_n \ge 0\), this means \(y_n (\bm{w}^\top
                \bm{z}_n + b) \le 1\): the example is on the edge, inside the
                strip, or on the wrong side. Which of these it is cannot be
                read from \(\alpha_n\); we need \(b\) to compute \(\xi_n\),
                and then the list in \autoref{sec:soft-geometry} tells us
                where the example lies.
        \end{itemize}
\end{itemize}

These conclusions go in one direction only: from \(\alpha_n\) to the
position. The reverse does not always work. An example exactly on the edge
of the strip can have \(\alpha_n = 0\), any value in \((0, C)\), or \(\alpha_n
= C\), because all three cases allow \(y_n (\bm{w}^\top \bm{z}_n + b) = 1\).
That is what the greyed-out `\gone{on} fat boundary' in the first and third
lines of the slide refers to.

\begin{note}
    The fifteen points are measured from the slide, and the SVM (with \(C =
    100\)) is solved by us. The squares and triangles come from our
    solution. They agree with the deck except for two crosses at the left
    edge of the strip, which are almost equally far from the line, so the
    measurement cannot tell which of them lies exactly on the edge.
\end{note}

\begin{moral}
    \(\dvar{\alpha_n}\) can be used for \hlr{\textbf{data analysis}}.
\end{moral}

For example, the bounded SVs are the examples the SVM gave up on, or could
only partly fit. If we suspect that some labels are wrong, they are the
first examples to check.

\begin{exercise}[Fun Time]
    For a data set of size \(10000\), after solving SVM, assume that there
    are \(1126\) support vectors, and \(1000\) of those support vectors are
    bounded. What is the possible range of \(\Ein(g_{\mathrm{SVM}})\) in
    terms of \(0/1\) error?
    \begin{enumerate}[(1)]
        \item \(0.0000 \le \Ein(g_{\mathrm{SVM}}) \le 0.1000\);
        \item \(0.1000 \le \Ein(g_{\mathrm{SVM}}) \le 0.1126\);
        \item \(0.1126 \le \Ein(g_{\mathrm{SVM}}) \le 0.5000\);
        \item \(0.1126 \le \Ein(g_{\mathrm{SVM}}) \le 1.0000\).
    \end{enumerate}
\end{exercise}

\begin{proof}[Reference Answer]
    (1). The bounded support vectors are the only ones that could violate the
    fat boundary: \(\xi_n \ge 0\). If \(\xi_n \ge 1\), then the violation
    causes a \(0/1\) error on the example. On the other hand, it is also
    possible that \(\xi_n < 1\), and in that case the violation does not
    cause a \(0/1\) error. So anywhere from none to all \(1000\) of them can
    be errors, which gives \(0 \le \Ein \le 1000/10000\).
\end{proof}

\section{Model Selection}\label{sec:soft-model-selection}

\source{slides p.17--21}

The soft-margin Gaussian SVM has two parameters, \(C\) and \(\gamma\), and
we have seen that both can cause overfitting. How should we choose them? This
is the model selection problem of \autoref{sec:model-selection}, and the
standard answer is validation. At the end of the section we will also find a
cheap shortcut that works only for the SVM.

\subsection{Practical Need: Model Selection}

The nine pictures below train the soft-margin Gaussian SVM on the same noisy
data, one picture for each combination of \(C\) and \(\gamma\). The columns
go from \(C = 1\) to \(C = 100\); the rows go from large \(\gamma\) (top) to
small \(\gamma\) (bottom). The middle row is the three pictures from the
previous section.

\begin{center}
    \begin{minipage}[c]{0.56\textwidth}
        \centering
        \smsvmgrid{}{}{}{}{}{}{}{}{}
    \end{minipage}%
    \hfill
    \begin{minipage}[c]{0.4\textwidth}
        \begin{itemize}
            \item complicated even for \(\hlo{(C, \gamma)}\)
                \hlo{\textbf{of Gaussian SVM}}
            \item more combinations if including other kernels or
                parameters
        \end{itemize}
    \end{minipage}
\end{center}

\begin{note}
    The slides do not give the values of \(\gamma\). The values shown are in
    our own coordinates (the box is \([-1, 1]^2\)), chosen so that each
    picture matches the deck's.
\end{note}

In the top row (large \(\gamma\)) the boundary winds around almost every
example, whatever \(C\) is. The bottom-left picture (small \(\gamma\), small
\(C\)) has the smoothest boundary. But which one will predict best on
\emph{new} data? We cannot tell by looking; \autoref{sec:visual-choices}
explained why choosing by eye is dangerous. With other kernels there would be
even more to choose, such as \((\gamma, \zeta, Q)\) for the polynomial
kernel.

\begin{moral}
    How to select? \hlo{\textbf{Validation :-)}}
\end{moral}

\subsection{Selection by Cross Validation}

Treat each combination \((C, \gamma)\) as one model. \(V\)-fold cross
validation (\autoref{sec:vfold}) gives each model a score \(\Ecv(C,
\gamma)\), an estimate of its \(\Eout\). The numbers on the panels are the
deck's \(\Ecv\) values.

\begin{center}
    \begin{minipage}[c]{0.56\textwidth}
        \centering
        \smsvmgrid{0.3500}{0.3250}{0.3250}{0.2000}{0.2250}{0.2750}%
            {\hlo{\textbf{0.1750}}}{0.2250}{0.2000}
    \end{minipage}%
    \hfill
    \begin{minipage}[c]{0.4\textwidth}
        \begin{itemize}
            \item \(\Ecv(C, \gamma)\): `non-smooth' function of \((C,
                \gamma)\) --- \hlo{\textbf{difficult to optimize}}
            \item proper models can be chosen by \hlo{\textbf{\(V\)-fold
                cross validation}} on \hlo{\textbf{a few grid values of
                \((C, \gamma)\)}}
        \end{itemize}
    \end{minipage}
\end{center}

Why is \(\Ecv\) `non-smooth'? It is an average of \(0/1\) errors, so when
\((C, \gamma)\) changes a little, \(\Ecv\) either stays the same or jumps.
There is no gradient to follow, so we cannot optimize it the way we optimize
a smooth error. The practical method is a grid search: try a few values of
each parameter (often a factor of \(10\) apart, as here), compute \(\Ecv\)
for every combination, and keep the best one. Here that is the bottom-left
model, with \(\Ecv = 0.1750\).

\begin{moral}
    \(\Ecv\): very popular criteria for soft-margin SVM.
\end{moral}

\subsection{Leave-One-Out CV Error for SVM}

\begin{recall}
    \(\Eloocv = \Ecv\) with \(N\) folds (\autoref{sec:loocv}).
\end{recall}

Leave-one-out CV is the most expensive kind of cross validation: for each
model it trains \(N\) times, once with each example left out. For the SVM,
however, most of those trainings are unnecessary, and this gives a cheap
upper bound on \(\Eloocv\).

\begin{center}
    \begin{minipage}[c]{0.56\textwidth}
        \begin{theorem}[Claim]\label{thm:loocv-sv}
            \[
                \Eloocv \le \frac{\#\text{SV}}{N}.
            \]
        \end{theorem}
    \end{minipage}%
    \hfill
    \begin{minipage}[c]{0.4\textwidth}
        \centering
        \svmtoysolution\\[2pt]
        {\small motivation from hard-margin SVM: only \hlo{\textbf{SVs
        needed}}}
    \end{minipage}
\end{center}

The idea comes from the toy example of \autoref{sec:svm}, shown on the
right. Only the three boxed examples (the SVs) determine the fattest
hyperplane. If we remove the fourth example and train again, we get exactly
the same hyperplane, and it still classifies the removed example correctly.
So that example contributes \(0\) to \(\Eloocv\). The proof below shows that
this holds for every non-SV.

\begin{proof}
    Look at the example \(\hlb{(\bm{x}_N, y_N)}\); any other example works
    the same way after renumbering. Write \(F(\alpha_1, \dots, \alpha_N) =
    \frac{1}{2} \sum_n \sum_m \alpha_n \alpha_m q_{n,m} - \sum_n \alpha_n\)
    for the dual objective. Training on all \(N\) examples minimizes \(F\)
    over \(\alpha_1, \dots, \alpha_N\) subject to \(\sum_{n=1}^N y_n \alpha_n
    = 0\) and \(0 \le \alpha_n \le C\). Training with \((\bm{x}_N, y_N)\) left
    out is the same problem without \(\alpha_N\): minimize \(F(\alpha_1,
    \dots, \alpha_{N-1}, 0)\) subject to \(\sum_{n=1}^{N-1} y_n \alpha_n = 0\)
    and \(0 \le \alpha_n \le C\).
    \begin{itemize}
        \item Suppose its optimal \(\dvar{\alpha_N = 0}\) (non-SV). We claim
            that \(\dvar{(\alpha_1, \alpha_2, \dots , \alpha_{N-1})}\) is
            \dvar{still optimal} when \hlb{leaving out \((\bm{x}_N,
            y_N)\)}, i.e.\ it solves the smaller problem.

            First, it satisfies the constraints of the smaller problem:
            \(\sum_{n<N} y_n \alpha_n = \sum_n y_n \alpha_n = 0\) because
            \(\alpha_N = 0\), and each \(\alpha_n\) is still in \([0, C]\).
            Second, its objective value is the same as that of the full
            \(\bm{\alpha}\), because every term with \(\alpha_N\) is zero.

            The key: \textbf{what if there's better} \(\alpha_n\) for the
            smaller problem? Add \(\alpha_N = 0\) back to it as the
            \(N\)-th entry. The result satisfies the constraints of the full
            problem and has the same, smaller value of \(F\). It would beat
            the optimal \(\bm{\alpha}\) of the full problem, which is
            impossible.
        \item So for the SVM, \(\hlb{g^-} = g\) when \hlb{leaving out
            non-SV}: the same \(\alpha_n\) give the same \(\bm{w} = \sum_n
            \alpha_n y_n \bm{z}_n\), and the same \(b\).
        \item Therefore
            \[
                e_{\text{non-SV}} = \err(\hlb{g^-}, \hlb{\text{non-SV}}) =
                \err(g, \hlb{\text{non-SV}}) = 0.
            \]
            The last error is \(0\) because a non-SV has \(\xi_n = 0\), so
            \(y_n (\bm{w}^\top \bm{z}_n + b) \ge 1\): \(g\) classifies it
            correctly.
        \item For an SV we only know \(e_{\hlo{\text{SV}}} \le 1\).
    \end{itemize}
    Adding up, only the SVs can contribute errors, and each contributes at
    most \(1\): \(\Eloocv = \frac{1}{N} \sum_n e_n \le
    \frac{\#\text{SV}}{N}\).
\end{proof}

\begin{remark}
    The proof says `the same \(\bm{w}\) and the same \(b\)'. But the smaller
    problem may have several optimal solutions. Is \(g^-\) then really
    \(g\)?
    \begin{itemize}
        \item \(\bm{w}\) is always the same. The optimal \(\bm{w}\) of the
            primal is unique, because \(\frac{1}{2} \bm{w}^\top \bm{w}\) is
            strictly convex, and every optimal \(\bm{\alpha}\) gives this
            \(\bm{w}\) through \(\bm{w} = \sum_n \alpha_n y_n \bm{z}_n\).
        \item \(b\) is the same if \(g\) has a free SV. That example is
            still a free SV in the smaller problem, since its
            \(\alpha_n\) did not change, and it fixes \(b\) by
            \autoref{sec:soft-solve-b}.
        \item Without a free SV, \(b\) is only known to lie in an interval,
            and removing a non-SV can make this interval wider. The \(b\)
            used by \(g\) is still optimal for the smaller problem, but the
            solver may pick another one. Then \(g^- \ne g\), and \(g^-\) may
            misclassify the removed example, so \(e_{\text{non-SV}} = 0\)
            is no longer guaranteed.
    \end{itemize}
    So the bound holds when the solver picks the same \(b\) in both
    trainings. Since free SVs almost always exist, this is rarely a concern
    in practice.
\end{remark}

\begin{moral}
    Scaled \#SV bounds leave-one-out CV error.
\end{moral}

\subsection{Selection by \texorpdfstring{\#}{\#}SV}

Cross validation is expensive: \(V\)-fold CV trains every model \(V\) times,
and leave-one-out CV trains it \(N\) times. \autoref{thm:loocv-sv} offers a
cheaper quantity.

\begin{definition}
    For parameters \((C, \gamma)\), train the soft-margin SVM once on all
    \(N\) examples. Its \vocab{number of support vectors}
    \(n_{\mathrm{SV}}(C, \gamma)\) is the number of examples with
    \(\alpha_n > 0\) in the solution.
\end{definition}

By \autoref{thm:loocv-sv}, \(\Eloocv(C, \gamma) \le n_{\mathrm{SV}}(C,
\gamma) / N\). So one training per model gives an upper bound on its
leave-one-out error. The idea of this section is to use \(n_{\mathrm{SV}}\)
in place of \(\Ecv\) when \(\Ecv\) is too costly. The numbers on the panels
are the deck's \(n_{\mathrm{SV}}\) values for the same nine models as
before, with \(N = 40\).

\begin{center}
    \begin{minipage}[c]{0.56\textwidth}
        \centering
        \smsvmgrid{38}{37}{37}{27}{21}{\hlo{\textbf{17}}}%
            {21}{18}{19}
    \end{minipage}%
    \hfill
    \begin{minipage}[c]{0.4\textwidth}
        \begin{itemize}
            \item \(n_{\mathrm{SV}}(C, \gamma)\): `non-smooth' function of
                \((C, \gamma)\) --- \hlo{\textbf{difficult to optimize}}
            \item \hlo{\textbf{just an upper bound!}}
            \item dangerous models can be ruled out by
                \hlo{\(\bm{n_{\mathrm{SV}}}\)} on \hlo{\textbf{a few grid
                values of} \((C, \gamma)\)}
        \end{itemize}
    \end{minipage}
\end{center}

The three bullets say how \(n_{\mathrm{SV}}\) can and cannot be used.
\begin{itemize}
    \item \textbf{It is non-smooth.} \(n_{\mathrm{SV}}\) is a count, so it
        jumps as \((C, \gamma)\) changes, just like \(\Ecv\). We again
        evaluate it on a small grid instead of optimizing it.
    \item \textbf{It is only an upper bound.} A small \(n_{\mathrm{SV}}\)
        means a small bound, not a small error. So we should not pick the
        model with the smallest \(n_{\mathrm{SV}}\). Here the smallest
        count, \(17\) at \(\gamma = 5\), \(C = 100\), belongs to the model
        whose \(\Ecv = 0.2750\) is the worst in the bottom two rows.
    \item \textbf{It can rule out bad models.} A large \(n_{\mathrm{SV}}\)
        is a warning. In the top row, \(37\) or \(38\) of the \(40\)
        examples are SVs. The bound \(38/40 = 0.95\) says nothing useful,
        and a model that needs almost every example as an SV is memorizing
        the data; its boundary winds around each example. We can discard
        these models without running any CV.
\end{itemize}
In short, \(n_{\mathrm{SV}}\) filters out dangerous models cheaply; choosing
among the remaining ones is still the job of \(\Ecv\), if we can afford it.

\begin{note}
    Our own solutions give slightly different counts: \(39, 39, 39\) / \(28,
    24, 20\) / \(23, 19, 17\). The small differences come from the unknown
    coordinates of the deck and from the solver's tolerance.
\end{note}

\begin{moral}
    \(n_{\mathrm{SV}}\): often used as a \hlo{\textbf{safety check}} if computing
    \(\Ecv\) is too time-consuming.
\end{moral}

\begin{exercise}[Fun Time]
    For a data set of size \(10000\), after solving SVM on some parameters,
    assume that there are \(1126\) support vectors, and \(1000\) of those
    support vectors are bounded. Which of the following \emph{cannot} be
    \(\Eloocv\) with those parameters?
    \begin{enumerate}[(1)]
        \item \(0.0000\);
        \item \(0.0805\);
        \item \(0.1111\);
        \item \(0.5566\).
    \end{enumerate}
\end{exercise}

\begin{proof}[Reference Answer]
    (4). Note that the upper bound of \(\Eloocv\) is \(1126 / 10000 =
    0.1126\) by \autoref{thm:loocv-sv}.
\end{proof}

\hrulebar

\subsection*{Summary}

\begin{description}
    \item[Motivation and Primal Problem] \hlo{\textbf{add margin violations
        \(\bm{\xi_n}\)}}.
    \item[Dual Problem] \hlo{\textbf{upper-bound \(\bm{\alpha_n}\) by
        \(\bm{C}\)}}.
    \item[Messages behind Soft-Margin SVM] \hlo{\textbf{bounded/free SVs for
        data analysis}}.
    \item[Model Selection] \hlo{\textbf{cross-validation, or approximately
        \(\bm{n_{\mathrm{SV}}}\)}}.
\end{description}

The soft-margin SVM lets examples violate the margin and pays \(C \xi_n\)
for each violation, where \(\xi_n = \max\bigl(0, 1 - y_n (\bm{w}^\top
\bm{z}_n + b)\bigr)\). This looks a lot like an error measure plus a
regularizer, the form we met in regularization. The next chapter follows
that idea. It reads the soft-margin SVM as a regularized model, compares it
with logistic regression, and builds a kernel version of logistic
regression. \hlo{\textbf{Next: other kernel models for soft binary
classification.}}
